\chapter{Glossary}\label{app:glossary}

\begin{description}[leftmargin=3.2cm,style=nextline]
\item[Amplitude] the size of a signal; in RS3D in \emph{counts} or, after conversion, in m/s.
\item[Band-pass filter] keeps the frequencies between two limits and removes the rest.
\item[Channel] one stream of samples: \code{EHZ} vertical, \code{EHN} north--south, \code{EHE} east--west.
\item[Coda] the gradually fading tail of seismic waves after the main arrivals of an earthquake.
\item[Counts] the integer numbers produced by the Shake's analogue-to-digital converter; not a physical unit.
\item[Gap] time for which no samples were recorded.
\item[Helicorder] a plot that wraps a long record into rows, like lines of text.
\item[Hertz (Hz)] cycles per second, the unit of frequency.
\item[Link-local address] an address of the form \code{169.254.x.x} that devices give themselves when no router is present.
\item[miniSEED] the standard file format for seismic waveform data.
\item[Noise spectrum] the typical level of each frequency in a selected time range.
\item[P wave, S wave] the first (compressional) and second (shear) arrivals of an earthquake; the S wave arrives later and is
  usually stronger.
\item[SDS] SeisComP data structure: one file per channel per day in a fixed folder layout.
\item[SeedLink] the protocol the Shake uses to send data live over the network.
\item[Spectrogram] a picture of frequency against time, coloured by power.
\item[SSH] a secure login used by RS3D to read the Shake's clock, storage and archive.
\item[UTC] coordinated universal time, the time scale used for all seismic data.
\end{description}
